% ============================================================
%  Configuration as Operational Law:
%  A Mathematical Foundation for Bounded Admission
%  in Human and Machine Systems
%
%  Sean Chatman
%  2026
% ============================================================

\documentclass[12pt, letterpaper]{article}

% ---- Geometry -----------------------------------------------
\usepackage[
  top=1.25in, bottom=1.25in,
  left=1.25in, right=1.25in
]{geometry}

% ---- Math ---------------------------------------------------
\usepackage{amsmath, amsthm, amssymb, amsfonts}
\usepackage{mathtools}
\usepackage{bm}

% ---- Typography ---------------------------------------------
\usepackage{microtype}
\usepackage{lmodern}
\usepackage[T1]{fontenc}
\usepackage[utf8]{inputenc}
\usepackage{setspace}
\onehalfspacing

% ---- Graphics and Tables ------------------------------------
\usepackage{booktabs}
\usepackage{array}
\usepackage{tikz}
\usetikzlibrary{arrows.meta, positioning, shapes.geometric, fit, calc}

% ---- References and Links -----------------------------------
\usepackage{hyperref}
\hypersetup{
  colorlinks=true,
  linkcolor=black,
  citecolor=black,
  urlcolor=black
}

% ---- Listings -----------------------------------------------
\usepackage{listings}
\lstset{
  basicstyle=\ttfamily\small,
  breaklines=true,
  frame=single,
  framesep=4pt,
  xleftmargin=12pt,
  xrightmargin=12pt,
}

% ---- Theorem Environments -----------------------------------
\theoremstyle{plain}
\newtheorem{theorem}{Theorem}[section]
\newtheorem{lemma}[theorem]{Lemma}
\newtheorem{proposition}[theorem]{Proposition}
\newtheorem{corollary}[theorem]{Corollary}

\theoremstyle{definition}
\newtheorem{definition}[theorem]{Definition}
\newtheorem{example}[theorem]{Example}
\newtheorem{axiom}[theorem]{Axiom}

\theoremstyle{remark}
\newtheorem{remark}[theorem]{Remark}

% ---- Operators ----------------------------------------------
\DeclareMathOperator{\Ker}{Ker}
\DeclareMathOperator{\Img}{Im}
\DeclareMathOperator{\id}{id}
\DeclareMathOperator{\dom}{dom}
\DeclareMathOperator{\cod}{cod}
\DeclareMathOperator{\rank}{rank}
\DeclareMathOperator{\tr}{tr}
\DeclareMathOperator{\vol}{vol}
\DeclareMathOperator{\conv}{conv}
\DeclareMathOperator{\dist}{dist}

% ---- Custom Notation ----------------------------------------
\newcommand{\R}{\mathbb{R}}
\newcommand{\N}{\mathbb{N}}
\newcommand{\Z}{\mathbb{Z}}
\newcommand{\Q}{\mathbb{Q}}
\newcommand{\B}{\mathbb{B}}   % Boolean
\newcommand{\C}{\mathcal{C}}  % Configuration space
\newcommand{\A}{\mathcal{A}}  % Admitted set
\newcommand{\F}{\mathcal{F}}  % Field
\newcommand{\Lc}{\mathcal{L}} % Layer lattice
\newcommand{\W}{\mathcal{W}}  % Witness algebra
\newcommand{\ce}{\mathcal{E}} % Counterexample set
\newcommand{\qc}{q_{\mathrm{config}}}
\newcommand{\Phi}{\Phi}
\newcommand{\pword}[1]{\textsc{#1}}
\newcommand{\norm}[1]{\left\lVert #1 \right\rVert}
\newcommand{\abs}[1]{\left\lvert #1 \right\rvert}
\newcommand{\inner}[2]{\left\langle #1,\, #2 \right\rangle}
\newcommand{\floor}[1]{\left\lfloor #1 \right\rfloor}
\newcommand{\ceil}[1]{\left\lceil #1 \right\rceil}
\newcommand{\bigland}{\bigwedge}

% ============================================================
\begin{document}

% ---- Title Page --------------------------------------------
\begin{titlepage}
\centering
\vspace*{2cm}

{\LARGE\bfseries Configuration as Operational Law}\\[0.6em]
{\large A Mathematical Foundation for Bounded Admission\\
in Human and Machine Systems}

\vspace{2.5cm}

{\large Sean Chatman}\\[0.4em]
{\normalsize Pasadena, California}\\[0.2em]
{\normalsize \texttt{sean@chatmangpt.com}}

\vspace{2cm}

{\normalsize June 2026}

\vspace{3cm}

\begin{abstract}
\noindent
Configuration is the first operational law of any system that acts on instructions.
Every field value is a bounded assertion; every merge is a structured judgment;
every deployment is an execution of law that either has standing or does not.
This paper introduces \emph{config admission} as a formal discipline grounded
in set theory, lattice geometry, abstract algebra, real analysis, and
measure-theoretic process intelligence.
The central standing criterion---the \emph{Bounded Receipted Chatman Equation}---is
defined as a conjunction of four checkable predicates over a configuration
space $\C$, and proved to be sound, monotone, and computable.
A typestate admission pipeline is characterized as a monoidal category
whose terminal object is the cryptographically witnessed admitted artifact.
The derivative of standing with respect to configuration perturbations is
derived to support sensitivity analysis and authoring-time feedback.
Design for Combinatorial Maximalism (DfCM) is formalized as a
measure over the variant space, and OCEL telemetry is characterized as a
filtered probability space that captures process history without constituting
standing authority.
The paper closes with a generalization: bounded admission is the correct
model for any artifact---configuration, agent output, build plan,
deployment manifest, or policy file---that instructs a machine or
institution to act.
The future is not autonomous generation.
The future is bounded admission of generated artifacts.
\end{abstract}

\vfill
\end{titlepage}

% ---- Table of Contents -------------------------------------
\tableofcontents
\newpage

% ============================================================
\section{Prologue: The Hidden Operating System}

Every system that acts on instructions has a configuration layer.
In software, this layer appears as TOML, YAML, JSON, environment variables,
feature flags, infrastructure declarations, and runtime profiles.
In human institutions, it appears as forms, rubrics, schedules,
eligibility rules, care plans, work orders, and operating procedures.
In machine learning systems, it appears as hyperparameter files,
agent instructions, prompt templates, and tool registries.

These are all the same class of object: a bounded assertion layer
that governs what may happen before execution begins.

For decades, this layer has been treated as text---read by humans,
modified by convention, certified by prose.
The consequences are well-documented.
Systems fail not because they lack instructions,
but because their instructions are implicit, stale, conflicting, unbounded,
unwitnessed, or prose-certified.

\begin{definition}[Operational Law]
An \emph{operational law} for a system $S$ is a structured artifact
$\ell \in \C$ drawn from a configuration space $\C$,
together with a predicate $\qc : \C \to \B$ such that
$S$ may execute under $\ell$ only if $\qc(\ell) = 1$.
\end{definition}

This paper constructs $\qc$ from first principles.

% ============================================================
\section{Configuration Space: Geometry and Topology}

\subsection{The Configuration Space}

Let $\F$ be a value domain---a structured collection of types including
booleans, integers, reals, strings, enumerations, structured records,
sequences, and maps.
A \emph{configuration space} $\C$ is a product space over a finite
schema $\Sigma$:

\begin{equation}
  \C = \prod_{k \in \Sigma} \F_k
\end{equation}

where each $\F_k$ is the value domain of field $k$.
Points $x \in \C$ are configurations: total assignments of values to fields.
Partial configurations form a subspace $\C_\partial \subset \C^\bot$
where $\bot$ denotes the undefined sentinel.

\subsection{Constraint Sets and the Feasible Region}

Each field $k \in \Sigma$ admits a \emph{constraint function}
$\phi_k : \F_k \to \B$.
The set of constraint functions induces a \emph{feasible region}:

\begin{equation}
  \A_0 = \left\{ x \in \C \;\middle|\; \forall k \in \Sigma,\; \phi_k(x_k) = 1 \right\}
\end{equation}

When each $\F_k$ carries a natural ordering (e.g., numeric ranges),
$\A_0$ is the intersection of finitely many halfspaces and is therefore
a convex polytope in the continuous case.

\begin{proposition}[Feasibility as Intersection]
Let $H_k = \phi_k^{-1}(1) \subseteq \F_k$ be the feasible set for field $k$.
Then $\A_0 = \prod_{k \in \Sigma} H_k$.
When each $H_k$ is convex, $\A_0$ is convex.
\end{proposition}

\begin{proof}
Immediate from the definition of the product topology and the
fact that finite intersections of convex sets are convex.
\end{proof}

\begin{definition}[Constraint Boundary]
The \emph{constraint boundary} of $\A_0$ is
$\partial \A_0 = \overline{\A_0} \setminus \mathrm{int}(\A_0)$.
Points on $\partial \A_0$ satisfy all constraints with equality in at
least one coordinate.
They are valid but maximally constrained.
\end{definition}

\begin{remark}
For discrete value domains (enumerations, boolean fields),
$\A_0$ is a finite set and the geometric language of convexity
applies via the convex hull: $\conv(\A_0) \subseteq \R^n$
embeds $\A_0$ in a real vector space by coordinate-encoding.
\end{remark}

\subsection{The Geometry of Violation}

Define the \emph{violation set}:
\begin{equation}
  V = \C \setminus \A_0 = \left\{ x \in \C \;\middle|\;
  \exists k \in \Sigma,\; \phi_k(x_k) = 0 \right\}
\end{equation}

The violation set is the complement of the feasible region.
Its geometry determines the difficulty of error correction:
configurations near $\partial \A_0$ require small repairs;
configurations deep in $V$ require large ones.

\begin{definition}[Violation Distance]
For $x \in V$, the \emph{violation distance} is:
\begin{equation}
  d(x, \A_0) = \inf_{y \in \A_0} \norm{x - y}
\end{equation}
For a product metric $\norm{x - y}^2 = \sum_k w_k (x_k - y_k)^2$
with field weights $w_k > 0$, this is the weighted Euclidean distance
to the nearest valid configuration.
\end{definition}

The violation distance is the natural quantity for authoring-time feedback:
it tells the operator how far the current configuration is from standing.

\subsection{Layer Geometry: The Merge Polytope}

Modern configuration systems compose values from multiple sources
(files, environment variables, command-line overrides) using a
\emph{layer precedence order}.
Let $[n] = \{1, \ldots, n\}$ index the layers in increasing precedence.
Each layer $i$ contributes a partial configuration $x^{(i)} \in \C_\partial$.

The \emph{merged configuration} under winner-takes-all precedence is:

\begin{equation}
  x^* = \bigsqcup_{i=1}^{n} x^{(i)}
\end{equation}

where $\sqcup$ denotes the \emph{precedence join}: for each field $k$,
$x^*_k = x^{(i^*)}_k$ where $i^* = \max\{i \mid x^{(i)}_k \neq \bot\}$.

\begin{definition}[Winner Map]
The \emph{winner map} $W : \Sigma \to [n] \cup \{\bot\}$ assigns
to each field the index of the layer whose value prevails:
\begin{equation}
  W(k) = \max\left\{ i \in [n] \;\middle|\; x^{(i)}_k \neq \bot \right\}
\end{equation}
$W(k) = \bot$ if no layer provides a value for field $k$.
\end{definition}

The winner map is the \emph{provenance certificate} for the merge.
A merge is \emph{lawful} if and only if $W$ is computable, deterministic,
and total over the required schema.

\begin{theorem}[Merge Uniqueness]
For a fixed layer ordering $\sigma \in S_n$ (a permutation of $[n]$),
the merged configuration $x^*_\sigma$ is unique.
Permuting $\sigma$ produces (in general) a different merge.
\end{theorem}

\begin{proof}
For each $k$, $W_\sigma(k) = \max_\sigma\{i \mid x^{(i)}_k \neq \bot\}$
is determined entirely by the ordering $\sigma$ and the support sets
$\{i \mid x^{(i)}_k \neq \bot\}$.
Since $\max_\sigma$ is well-defined for any total order $\sigma$,
$x^*_\sigma$ is unique.
Reordering $\sigma$ changes which layer wins contested fields,
so in general $x^*_\sigma \neq x^*_{\sigma'}$ for $\sigma \neq \sigma'$.
\end{proof}

\begin{corollary}
A merge is \emph{order-independent} (commutative) if and only if
no field is defined in more than one layer.
In all other cases, the layer ordering is a required part of the
configuration artifact.
\end{corollary}

% ============================================================
\section{Algebraic Structure of Admission}

\subsection{The Admission Lattice}

\begin{definition}[Admission Lattice]
The \emph{admission lattice} $(\Lc, \leq)$ is the power set of $\Sigma$
ordered by inclusion, where a set $S \subseteq \Sigma$ of validated fields
forms a \emph{partial admission state}.
The bottom element $\hat{0} = \emptyset$ represents no validated fields;
the top element $\hat{1} = \Sigma$ represents full admission.
\end{definition}

The lattice $(\Lc, \leq, \wedge, \vee)$ is a complete Boolean lattice
isomorphic to $2^{|\Sigma|}$.
It is the algebraic structure underlying the incremental validation
process: each check extends the partial admission state monotonically.

\begin{proposition}[Monotonicity of Validation]
If $S \leq S'$ in $\Lc$ (i.e., $S \subseteq S'$), then
any field-level constraint $\phi_k$ that holds for configuration
$x$ at state $S'$ continues to hold at $S$.
Validation progress is monotone: checks do not un-pass.
\end{proposition}

\begin{proof}
Field constraints $\phi_k$ are functions of individual field values,
not of validation state.
Passing a constraint at step $S$ means $\phi_k(x_k) = 1$,
which is independent of how many other fields have been checked.
\end{proof}

\subsection{The Typestate Monoid}

The admission pipeline is a sequence of transformations.
We formalize it as a \emph{typed monoid} acting on configuration states.

\begin{definition}[Admission Stage]
An \emph{admission stage} is a partial function
$f : \C_i \rightharpoonup \C_{i+1}$
between typed configuration spaces, where
$f$ is defined (i.e., the stage succeeds) if and only if
the input satisfies the stage's precondition.
\end{definition}

The full pipeline is the composition:

\begin{align}
  \pi = f_n \circ f_{n-1} \circ \cdots \circ f_1
\end{align}

where:
\begin{align*}
  f_1 &: \text{Raw} \to \text{BoundedSources} \\
  f_2 &: \text{BoundedSources} \to \text{Merged} \\
  f_3 &: \text{Merged} \to \text{EnvResolved} \\
  f_4 &: \text{EnvResolved} \to \text{Deserialized}(T) \\
  f_5 &: \text{Deserialized}(T) \to \text{Validated}(T) \\
  f_6 &: \text{Validated}(T) \to \text{Frozen}(T) \\
  f_7 &: \text{Frozen}(T) \to \text{CanonicalSaved} \\
  f_8 &: \text{CanonicalSaved} \to \text{Witnessed} \\
  f_9 &: \text{Witnessed} \to \text{AdmittedConfig}(T)
\end{align*}

\begin{theorem}[Pipeline as Monoidal Category]
The admission stages together with their typed composition form
a monoidal category $\mathbf{Adm}$ where:
\begin{enumerate}
  \item Objects are the stage types
        $\{\text{Raw}, \text{BoundedSources}, \ldots, \text{AdmittedConfig}(T)\}$.
  \item Morphisms are the stage functions $f_i$.
  \item Composition is function composition (when types match).
  \item The monoidal unit is $\id : \text{Raw} \to \text{Raw}$.
  \item There are no morphisms from \emph{later} stages to \emph{earlier}
        stages: the category is \emph{acyclic}.
\end{enumerate}
The terminal object in $\mathbf{Adm}$ is $\text{AdmittedConfig}(T)$.
\end{theorem}

\begin{proof}
Associativity of composition is inherited from function composition.
The identity morphism on each object is the identity function.
Acyclicity follows from the strict type discipline: each $f_i$ maps
between distinct types, and no inverse morphisms are defined.
The terminal property of $\text{AdmittedConfig}(T)$ holds because
no morphism exits it---it is the unique sink of the category.
\end{proof}

\begin{corollary}[Forgery Prevention]
Because $\text{AdmittedConfig}(T)$ is the terminal object of
$\mathbf{Adm}$ and no morphism constructs it outside of $\pi$,
an $\text{AdmittedConfig}(T)$ value cannot be constructed by
any path other than traversing the full pipeline.
\emph{Standing cannot be forged.}
\end{corollary}

\subsection{The Witness Group}

\begin{definition}[Witness Algebra]
Let $\W = (\{0,1\}^{256}, \oplus, \mathbf{0})$ be the group of
256-bit BLAKE3 digests under XOR, with identity element $\mathbf{0}$.
The \emph{witness function} $\omega : \A_0 \times \Lc^n \times E \to \W$
maps an admitted configuration $x$, layer ordering $\ell \in \Lc^n$,
and environment binding $e \in E$ to a witness value:
\begin{equation}
  \omega(x, \ell, e) = \mathrm{BLAKE3}\!\left(
    \bigoplus_{i} \mathrm{BLAKE3}(s_i)
    \;\|\;
    \mathrm{BLAKE3}(\ell)
    \;\|\;
    \mathrm{BLAKE3}(e)
    \;\|\;
    \mathrm{BLAKE3}(\kappa(x))
  \right)
\end{equation}
where $s_i$ are source file contents, $\|$ denotes concatenation,
and $\kappa(x)$ is the canonical serialization of $x$.
\end{definition}

\begin{theorem}[Witness Completeness and Determinism]
\label{thm:witness}
The witness function $\omega$ satisfies:
\begin{enumerate}
  \item \textbf{Determinism:} for fixed $(x, \ell, e)$,
        $\omega(x, \ell, e)$ is unique.
  \item \textbf{Sensitivity:} if any of $x$, $\ell$, or $e$ changes,
        then with overwhelming probability $\omega$ changes.
  \item \textbf{Preimage resistance:} given $\omega(x, \ell, e)$,
        it is computationally infeasible to recover $(x, \ell, e)$
        without knowing the inputs.
\end{enumerate}
\end{theorem}

\begin{proof}
(1) BLAKE3 is a deterministic function; all operations in $\omega$ are
deterministic given fixed inputs.
(2) Follows from the collision resistance of BLAKE3:
$\Pr[\omega(x,\ell,e) = \omega(x',\ell',e')] \leq 2^{-256}$
for $(x,\ell,e) \neq (x',\ell',e')$.
(3) Preimage resistance is a standard property of BLAKE3 as a
cryptographic hash function under the random oracle model.
\end{proof}

% ============================================================
\section{The Bounded Receipted Chatman Equation}

We now state the central standing criterion in full generality.

\subsection{The Four Predicates}

\begin{definition}[BoundSatisfied]
A configuration $x \in \C$ satisfies \emph{BoundSatisfied} if
every declared field constraint holds:
\begin{equation}
  \mathrm{BS}(x) = \bigwedge_{k \in \Sigma} \phi_k(x_k) = 1
  \iff x \in \A_0
\end{equation}
\end{definition}

\begin{definition}[TransformLawful]
A merge $x^* = \bigsqcup_\sigma x^{(i)}$ with winner map $W$ satisfies
\emph{TransformLawful} if the winner map is total over required fields
and computed by a deterministic, documented precedence rule $\sigma$:
\begin{equation}
  \mathrm{TL}(x^*, W, \sigma) =
  \left[\dom(W) \supseteq \Sigma_{\mathrm{req}}\right]
  \wedge
  \left[W \text{ is } \sigma\text{-computable}\right]
\end{equation}
where $\Sigma_{\mathrm{req}} \subseteq \Sigma$ is the set of required fields.
\end{definition}

\begin{definition}[WitnessComplete]
A configuration $x$ with layer ordering $\ell$ and environment $e$
satisfies \emph{WitnessComplete} if a valid witness $w = \omega(x, \ell, e)$
has been computed and stored:
\begin{equation}
  \mathrm{WC}(x, w) = \left[ w = \omega(x, \ell, e) \right] \wedge \left[ w \neq \mathbf{0} \right]
\end{equation}
\end{definition}

\begin{definition}[CounterexamplesAbsent]
Let $\ce = \{c_1, \ldots, c_m\}$ be a finite set of
\emph{counterexample detectors}, each a function
$c_i : \C \to \B$ where $c_i(x) = 1$ means detector $i$ fires.
A configuration $x$ satisfies \emph{CounterexamplesAbsent} if
no detector fires:
\begin{equation}
  \mathrm{CA}(x) = \bigwedge_{i=1}^{m} \lnot c_i(x) = 1
  \iff x \notin \bigcup_{i=1}^{m} c_i^{-1}(1)
\end{equation}
\end{definition}

\subsection{The Standing Bit}

\begin{definition}[Bounded Receipted Chatman Equation]
\label{def:brce}
For a configuration $x \in \C$, winner map $W$, layer ordering $\sigma$,
witness $w$, and counterexample set $\ce$, the \emph{standing bit} is:
\begin{equation}
  \boxed{
  \qc(x, W, \sigma, w, \ce) = 1
  \iff
  \mathrm{BS}(x) \wedge \mathrm{TL}(x, W, \sigma) \wedge \mathrm{WC}(x, w) \wedge \mathrm{CA}(x)
  }
\end{equation}
and $\qc = 0$ otherwise.
The set of configurations with standing is the \emph{admitted set}:
\begin{equation}
  \A = \{ x \in \C \mid \qc(x, W, \sigma, w, \ce) = 1 \}
\end{equation}
\end{definition}

\begin{theorem}[Soundness of BRCE]
The admitted set $\A$ satisfies:
\begin{enumerate}
  \item $\A \subseteq \A_0$ (admitted configurations are feasible)
  \item $\A \subseteq \C \setminus \bigcup_i c_i^{-1}(1)$ (no counterexample fires)
  \item Every $x \in \A$ carries a valid witness (receipted)
  \item The map $x \mapsto \qc(x)$ is computable in polynomial time
        in $|\Sigma| + m$
\end{enumerate}
\end{theorem}

\begin{proof}
(1) $\mathrm{BS}(x) = 1 \Rightarrow x \in \A_0$ by Definition 2.3.
(2) $\mathrm{CA}(x) = 1 \Rightarrow \forall i,\; c_i(x) = 0$ by Definition 4.4.
(3) $\mathrm{WC}(x,w) = 1 \Rightarrow w = \omega(x,\ell,e)$ by Definition 4.3.
(4) Each $\phi_k$ is computable in $O(1)$; their conjunction is $O(|\Sigma|)$.
$\mathrm{TL}$ checks $O(|\Sigma|)$ fields.
$\mathrm{WC}$ computes one BLAKE3 hash: $O(|x|)$.
$\mathrm{CA}$ evaluates $m$ detectors: $O(m)$.
Total: $O(|\Sigma| + m + |x|)$.
\end{proof}

\begin{theorem}[BRCE is Strictly Stronger than Parse Correctness]
Let $\mathrm{Parse}(x) = 1$ denote that $x$ is syntactically valid TOML.
Then $\qc(x) = 1 \Rightarrow \mathrm{Parse}(x) = 1$, but
$\mathrm{Parse}(x) = 1 \not\Rightarrow \qc(x) = 1$.
\end{theorem}

\begin{proof}
Forward direction: any $x \in \A$ was successfully deserialized at
stage $f_4$, which requires parse correctness.
Converse: take $x$ with a field value outside its declared range.
Then $\mathrm{BS}(x) = 0$, so $\qc(x) = 0$, but $x$ may parse correctly.
\end{proof}

\begin{corollary}[Prose Does Not Entail Standing]
Let $\Sigma_{\mathrm{prose}}$ be the set of natural-language assertions
about a configuration (README statements, review comments, approval emails).
Then for any $x$ and any prose set $P$:
\begin{equation}
  \Sigma_{\mathrm{prose}} \not\vdash \qc(x) = 1
\end{equation}
Standing is entailed only by the BRCE predicate, not by description.
\end{corollary}

% ============================================================
\section{Calculus of Standing}

We now develop a differential theory of configuration standing:
how does $\qc$ change as a configuration changes?
This theory grounds the authoring-time feedback loop.

\subsection{Smooth Extension of the Standing Predicate}

The standing bit $\qc : \C \to \{0,1\}$ is discontinuous.
To apply calculus, we construct a smooth surrogate
$\tilde{q} : \R^n \to [0,1]$ that agrees with $\qc$ on $\A$ and $V$.

For numeric fields with interval constraints $[a_k, b_k]$,
define the \emph{satisfaction score} for field $k$:

\begin{equation}
  s_k(x_k) = \exp\!\left( -\frac{[\max(a_k - x_k,\; x_k - b_k,\; 0)]^2}{2\sigma_k^2} \right)
\end{equation}

This is a Gaussian-envelope penalty: $s_k = 1$ when $x_k \in [a_k, b_k]$
and decays smoothly to $0$ as $x_k$ moves away from the feasible interval.
The parameter $\sigma_k > 0$ controls the rate of decay.

The \emph{smooth standing score} is:

\begin{equation}
  \tilde{q}(x) = \prod_{k \in \Sigma} s_k(x_k)
\end{equation}

Note that $\tilde{q}(x) = 1$ if and only if all fields are in their
feasible ranges, and $\tilde{q}(x) < 1$ otherwise.

\subsection{The Standing Gradient}

\begin{proposition}[Standing Gradient]
The gradient of $\tilde{q}$ with respect to the $k$-th field is:
\begin{equation}
  \frac{\partial \tilde{q}}{\partial x_k}
  = \tilde{q}(x) \cdot \frac{\partial \ln s_k}{\partial x_k}
\end{equation}
where
\begin{equation}
  \frac{\partial \ln s_k}{\partial x_k} =
  \begin{cases}
    \dfrac{x_k - a_k}{\sigma_k^2} & \text{if } x_k < a_k \\[6pt]
    0 & \text{if } x_k \in [a_k, b_k] \\[4pt]
    \dfrac{b_k - x_k}{\sigma_k^2} & \text{if } x_k > b_k
  \end{cases}
\end{equation}
\end{proposition}

\begin{proof}
By the chain rule: $\frac{\partial \tilde{q}}{\partial x_k}
= \tilde{q}(x) \cdot \frac{\partial \ln \tilde{q}}{\partial x_k}
= \tilde{q}(x) \cdot \frac{\partial \ln s_k}{\partial x_k}$,
since $\ln \tilde{q} = \sum_j \ln s_j$ and only the $k$-th term depends on $x_k$.
The formula for $\partial \ln s_k / \partial x_k$ follows from
differentiating the Gaussian envelope.
\end{proof}

The gradient points \emph{toward} the feasible region from outside it,
and is zero inside.
This is the mathematical basis for authoring-time repair hints:
the gradient $-\nabla \tilde{q}$ is a repair direction.

\subsection{Sensitivity Analysis}

\begin{definition}[Configuration Sensitivity]
The \emph{sensitivity} of the standing score to field $k$ is:
\begin{equation}
  \Lambda_k(x) = \left\lvert \frac{\partial \tilde{q}}{\partial x_k}(x) \right\rvert
\end{equation}
The \emph{sensitivity vector} $\bm{\Lambda}(x) = (\Lambda_1(x), \ldots, \Lambda_n(x))$
ranks fields by how much their perturbation affects standing.
\end{definition}

\begin{theorem}[Sensitivity Concentration]
For a product-form score $\tilde{q}(x) = \prod_k s_k(x_k)$:
\begin{equation}
  \Lambda_k(x) = \tilde{q}(x) \cdot \frac{\abs{r_k(x)}}{\sigma_k^2}
\end{equation}
where $r_k(x) = \max(a_k - x_k,\; x_k - b_k,\; 0)$ is the
\emph{constraint residual} for field $k$.
Sensitivity is concentrated on fields with large residuals
and small $\sigma_k$ (tight tolerances).
\end{theorem}

\begin{proof}
Follows directly from Proposition 5.2 and the definition of $r_k$.
\end{proof}

\subsection{The Repair Derivative}

Given a configuration $x \notin \A_0$ (lacking standing),
the \emph{steepest-ascent repair} is:
\begin{equation}
  \delta x^* = \arg\max_{\norm{\delta x} = \epsilon} \tilde{q}(x + \delta x)
  \approx \epsilon \cdot \frac{\nabla \tilde{q}(x)}{\norm{\nabla \tilde{q}(x)}}
\end{equation}
for small $\epsilon > 0$.
This is the direction of maximum improvement in the standing score
for a unit-norm perturbation.

\begin{theorem}[Convergence of Greedy Repair]
The gradient ascent iteration $x_{t+1} = x_t + \eta \nabla \tilde{q}(x_t)$
converges to a point in $\A_0$ for sufficiently small learning rate $\eta > 0$,
provided $\tilde{q}$ is strictly log-concave on the path from $x_0$ to $\A_0$.
\end{theorem}

\begin{proof}
$\tilde{q}$ is a product of Gaussian envelopes.
Each factor $s_k$ is log-concave (Gaussian densities are log-concave).
A product of log-concave functions is log-concave.
Gradient ascent on a smooth, bounded, log-concave objective converges
to a global maximum by standard results in convex optimization
(the log-concavity of $\tilde{q}$ implies the concavity of $\ln \tilde{q}$,
whose maximizer is in $\A_0$).
\end{proof}

This theorem is the mathematical foundation for the LSP repair hint:
a repair direction is computable, and iterative repair converges.

% ============================================================
\section{The Witness as a Receipt}

\subsection{Information Content of the Witness}

The witness $w = \omega(x, \ell, e)$ is a function of three inputs.
We characterize its information content using Shannon entropy.

\begin{definition}[Configuration Entropy]
Given a probability distribution $P$ over $\C$,
the \emph{configuration entropy} is:
\begin{equation}
  H(\C) = -\sum_{x \in \C} P(x) \log_2 P(x)
\end{equation}
The \emph{witness entropy} is $H(W) = H(\{0,1\}^{256}) = 256$ bits
(under the assumption that BLAKE3 acts as a random oracle).
\end{definition}

\begin{proposition}[Witness Compression Impossibility]
No two distinct admitted configurations $x \neq x'$ can share a witness.
\end{proposition}

\begin{proof}
Suppose $\omega(x, \ell, e) = \omega(x', \ell, e)$ with $x \neq x'$.
Since $x \neq x'$, their canonical serializations differ:
$\kappa(x) \neq \kappa(x')$.
Therefore $\mathrm{BLAKE3}(\kappa(x)) \neq \mathrm{BLAKE3}(\kappa(x'))$
(collision resistance of BLAKE3 with probability $1 - 2^{-256}$).
The outer hash then produces different witnesses.
Contradiction (with overwhelming probability).
\end{proof}

\subsection{The Witness as a Bijection}

\begin{theorem}[Witness Injectivity over Admitted Configurations]
The restriction $\omega|_{\A} : \A \to \W$ is injective
(with overwhelming cryptographic probability).
\end{theorem}

\begin{proof}
By Proposition 6.2 and Theorem 3.7 (collision resistance),
distinct admitted configurations produce distinct witnesses.
\end{proof}

This means the witness is a \emph{naming system} for admitted configurations:
every distinct configuration in $\A$ has a unique name in $\W$.
The name is not assigned by convention but computed from content.

% ============================================================
\section{Process Intelligence: Measure Theory of Standing}

\subsection{The Event Probability Space}

Configuration admission events unfold over time.
We model them as a filtered probability space.

\begin{definition}[Admission Probability Space]
Let $(\Omega, \mathcal{F}, P)$ be a probability space where:
\begin{itemize}
  \item $\Omega$ is the set of all possible configuration lifecycle trajectories
  \item $\mathcal{F} = (\mathcal{F}_t)_{t \geq 0}$ is a filtration representing
        the information available at time $t$
  \item $P$ is a probability measure over trajectories
\end{itemize}
A \emph{configuration event} is an $\mathcal{F}$-measurable random variable
$E : \Omega \to \mathcal{E}$ where $\mathcal{E}$ is the set of event types.
\end{definition}

\begin{definition}[OCEL as Filtered History]
The OCEL (Object-Centric Event Log) is the sequence of observed events
$\{E_{t_1}, E_{t_2}, \ldots, E_{t_N}\}$ up to time $T$,
which constitutes a realization of the stochastic process
$(E_t)_{t \geq 0}$ under the filtration $\mathcal{F}_T$.
OCEL is the history $\mathcal{F}_T$-measurable record.
It is \emph{not} a statement about $P(\qc = 1)$ at any future time.
\end{definition}

\begin{theorem}[OCEL is Not Standing Authority]
Let $\mathcal{H}_T = \sigma(E_{t_1}, \ldots, E_{t_N})$ be the
$\sigma$-algebra generated by all observed events up to time $T$.
Then $\mathcal{H}_T \not\vdash \qc(x) = 1$ for any $x$.
\end{theorem}

\begin{proof}
$\qc(x) = 1$ is a property of the \emph{current} configuration $x$
against the BRCE predicates (Definition \ref{def:brce}).
It depends on $\A_0$, $W$, $\sigma$, $\omega$, and $\ce$---none of which
are determined by the event history $\mathcal{H}_T$.
A record of past events cannot certify the current state of $x$
without re-running the admission pipeline on the current $x$.
\end{proof}

This theorem formalizes the architectural separation: OCEL is process history;
the witness is evidence; $\qc$ is the gate.
Conflating them is a category error.

\subsection{Design for Combinatorial Maximalism}

\begin{definition}[Variant Space]
The \emph{variant space} $\mathcal{V}$ is the Cartesian product of all
design axes for a configuration system:
\begin{equation}
  \mathcal{V} = \prod_{a \in \text{Axes}} D_a
\end{equation}
where $D_a$ is the domain of axis $a$.
Axes include source count, layer count, environment override count,
field type diversity, validation rule count, path policy, and
counterexample detector activation patterns.
\end{definition}

\begin{definition}[DfCM Measure]
The \emph{Design for Combinatorial Maximalism (DfCM) measure} is the
uniform probability measure $\mu$ over $\mathcal{V}$,
together with a target coverage function $\rho : 2^\mathcal{V} \to [0,1]$:
\begin{equation}
  \rho(S) = \frac{\mu(S)}{\mu(\mathcal{V})} = \frac{|S|}{|\mathcal{V}|}
\end{equation}
A test suite $T$ achieves \emph{combinatorial coverage} $\rho(T)$.
Full coverage means $\rho(T) = 1$; DfCM targets $\rho(T) \to 1$ as
resources allow.
\end{definition}

\begin{theorem}[Coverage Monotonicity]
For any test suites $T \subseteq T'$, $\rho(T) \leq \rho(T')$.
Adding tests never reduces combinatorial coverage.
\end{theorem}

\begin{proof}
$T \subseteq T' \Rightarrow |T| \leq |T'| \Rightarrow \rho(T) \leq \rho(T')$.
\end{proof}

\begin{definition}[Counterexample Density]
The \emph{counterexample density} of detector $c_i$ over the variant space is:
\begin{equation}
  \delta_i = \mu\!\left(\{ v \in \mathcal{V} \mid c_i(v) = 1 \}\right)
\end{equation}
A detector with high $\delta_i$ is a high-priority detector:
it fires across a large fraction of the variant space.
\end{definition}

The DfCM design discipline is to enumerate all axes,
measure the counterexample densities of all known detectors,
and prioritize coverage of high-density regions.

\subsection{The Fitness Score}

\begin{definition}[Fitness Score]
The \emph{fitness score} of a validation run is:
\begin{equation}
  f(x) = \frac{\abs{\{k \in \Sigma \mid \phi_k(x_k) = 1\}}}{\abs{\Sigma}}
  \in [0, 1]
\end{equation}
The fitness score measures the fraction of declared constraints that pass.
A configuration has standing only if $f(x) = 1$.
\end{definition}

\begin{proposition}[Fitness as Necessary but not Sufficient]
$\qc(x) = 1 \Rightarrow f(x) = 1$, but $f(x) = 1 \not\Rightarrow \qc(x) = 1$.
\end{proposition}

\begin{proof}
Forward: $\qc(x) = 1$ requires $\mathrm{BS}(x) = 1$, which means all
constraints pass, so $f(x) = 1$.
Converse: $f(x) = 1$ guarantees $\mathrm{BS}(x) = 1$ but not
$\mathrm{TL}$, $\mathrm{WC}$, or $\mathrm{CA}$.
A configuration with all field constraints passing but a missing witness
has $f = 1$ and $\qc = 0$.
\end{proof}

% ============================================================
\section{The Poka-Yoke Layer: Error Prevention Before Execution}

\subsection{Authoring-Time ANDON}

The admission pipeline runs at load time.
But defects introduced at authoring time---path traversal,
duplicate keys, unknown fields---should be caught earlier.
The language server provides this layer.

\begin{definition}[LSP Diagnostic]
An \emph{LSP diagnostic} is a tuple
$(r, \gamma, \delta, m)$ where:
\begin{itemize}
  \item $r = (\text{line}_1, \text{col}_1, \text{line}_2, \text{col}_2)$ is a source range
  \item $\gamma \in \{\text{Error}, \text{Warning}, \text{Information}\}$ is severity
  \item $\delta \in \Delta$ is a diagnostic code from the code set $\Delta$
  \item $m$ is a human-readable message
\end{itemize}
\end{definition}

\begin{definition}[ANDON Gate]
The \emph{ANDON gate} $\alpha : 2^\Delta \to \{0, 1\}$ maps the set of
active diagnostic codes to a binary signal:
\begin{equation}
  \alpha(D) =
  \begin{cases}
    1 & \text{if } D \cap \Delta_{\mathrm{error}} \neq \emptyset \\
    0 & \text{otherwise}
  \end{cases}
\end{equation}
where $\Delta_{\mathrm{error}} \subseteq \Delta$ is the set of error-severity codes.
$\alpha = 1$ means ANDON is raised (defect present);
$\alpha = 0$ means the file is diagnostically clean.
\end{definition}

\begin{theorem}[LSP is Not Standing Authority]
$\alpha(D) = 0 \not\Rightarrow \qc(x) = 1$.
\end{theorem}

\begin{proof}
$\alpha(D) = 0$ means no error-severity LSP diagnostics are active.
But $\qc(x) = 1$ requires $\mathrm{WC}$ (a witness must exist) and
$\mathrm{TL}$ (layer ordering must be verified)---properties that the
LSP, operating on a single file without full pipeline context,
cannot verify.
Therefore $\alpha(D) = 0$ is a necessary but not sufficient condition
for $\qc(x) = 1$.
\end{proof}

This is the correct architectural role of the LSP:
it is a \emph{poka-yoke} (mistake-proofing) device that prevents
common defects from reaching the admission pipeline,
not a substitute for the pipeline itself.

% ============================================================
\section{Generalization: Bounded Admission Beyond Configuration}

\subsection{The Admission Schema}

The BRCE framework applies to any artifact that:
\begin{enumerate}
  \item Has a declared schema with checkable field constraints
  \item Is composed from multiple sources under a precedence rule
  \item Can be canonicalized and witnessed
  \item Has a finite set of named failure detectors
\end{enumerate}

\begin{definition}[General Admission Problem]
Given an artifact space $\mathcal{X}$, a constraint family
$\Phi = \{\phi_k\}$, a merge operation $\sqcup$, a witness function $\omega$,
and a detector set $\ce$, the \emph{general admission problem} is to determine:
\begin{equation}
  q_{\mathcal{X}}(x) = \mathrm{BS}_\Phi(x)
    \wedge \mathrm{TL}_\sqcup(x)
    \wedge \mathrm{WC}_\omega(x)
    \wedge \mathrm{CA}_\ce(x)
\end{equation}
\end{definition}

The following table instantiates the four predicates for different artifact types:

\begin{table}[h]
\centering
\renewcommand{\arraystretch}{1.3}
\begin{tabular}{@{}lllll@{}}
\toprule
\textbf{Domain} & \textbf{BS} & \textbf{TL} & \textbf{WC} & \textbf{CA} \\
\midrule
Config file & Type/range & Layer order & BLAKE3 & Path traversal, dupes \\
Agent output & Schema match & Tool call order & Output hash & Hallucination detectors \\
Build plan & Dep bounds & Build order & Plan hash & Cycle detectors \\
Release desc & Version bounds & Changelog order & Release hash & Breaking-change detectors \\
Policy file & Rule bounds & Precedence & Policy hash & Contradiction detectors \\
Infra manifest & Resource bounds & Apply order & State hash & Drift detectors \\
\bottomrule
\end{tabular}
\caption{The four BRCE predicates instantiated across artifact domains.}
\end{table}

\subsection{The Phase Change}

The classical model of software deployment is:
\begin{equation}
  \text{generate} \to \text{test} \to \text{deploy} \to \text{observe}
\end{equation}
Configuration is assumed correct until observed to be wrong.

The admission model is:
\begin{equation}
  \text{generate} \to \text{bound} \to \text{merge} \to \text{witness} \to \text{admit} \to \text{execute}
\end{equation}
No artifact may execute without standing.

This is not an incremental improvement.
It is a phase change in the relationship between instruction and action.

\begin{definition}[Admitted Execution]
A system operates under \emph{admitted execution} if:
\begin{equation}
  \forall a \in \text{Actions},\; \text{precondition}(a)
  \supseteq \{\qc(\ell_a) = 1\}
\end{equation}
where $\ell_a$ is the operational law governing action $a$.
No action is taken without the corresponding law having standing.
\end{definition}

\begin{theorem}[Admitted Execution Reduces Unauthorized State]
Under admitted execution, the set of reachable system states is a
strict subset of the reachable states under unadmitted execution.
\end{theorem}

\begin{proof}
Let $R$ be the reachable state space under unadmitted execution and
$R_\A$ under admitted execution.
Under unadmitted execution, any configuration $\ell \in \C$ may
govern an action, including $\ell \notin \A$.
Under admitted execution, only $\ell \in \A$ governs actions.
Since $\A \subsetneq \C$ (the admitted set is a strict subset of all
configurations when any constraint is non-trivial),
$R_\A \subsetneq R$.
\end{proof}

% ============================================================
\section{The \texttt{star-toml} Implementation}

The foregoing theory is implemented in \texttt{star-toml},
a Rust library available at \texttt{https://github.com/seanchatmangpt/star-toml}.

\subsection{Typestate Realization}

The monoidal category $\mathbf{Adm}$ (Theorem 3.4) is realized in Rust
via \emph{typestate programming}: each stage type is a distinct Rust type
with private inner fields.
External code cannot construct the terminal object
\texttt{AdmittedConfig<T>} by struct literal---only the
\texttt{TrustedLoader::load\_admitted()} function produces it,
enforcing that the full pipeline has been traversed.

\subsection{BRCE Verification}

The \texttt{verifier\_report} binary checks 23 named counterexamples
(the set $\ce$ of Definition 4.4) against the current implementation.
Each check is a function $c_i : \text{State} \to \{0, 1\}$.
The verifier outputs a report with $|\ce_{\mathrm{active}}| = 0$ as
the standing condition for release.

\subsection{The LSP Realization}

The \texttt{star-toml-lsp} crate realizes the poka-yoke layer
(Section 8) as a Language Server Protocol server using \texttt{lsp-max}.
It computes $\alpha(D)$ on every document change and writes the ANDON gate
file \texttt{.star-toml/lsp.gate}.
It explicitly does not implement any path to \texttt{AdmittedConfig<T>}.

% ============================================================
\section{Conclusion}

We have derived the standing criterion $\qc$ from first principles:
\begin{itemize}
  \item \textbf{Geometry} gave us the feasible region $\A_0$,
        the violation distance, and the merge polytope.
  \item \textbf{Algebra} gave us the admission lattice $\Lc$,
        the monoidal category $\mathbf{Adm}$, and the witness group $\W$.
  \item \textbf{Calculus} gave us the standing gradient $\nabla \tilde{q}$,
        sensitivity analysis, and convergent repair.
  \item \textbf{Measure theory} gave us the DfCM coverage measure $\mu$,
        the fitness score $f$, and the formal separation of OCEL history
        from standing authority.
\end{itemize}

The four predicates---BoundSatisfied, TransformLawful, WitnessComplete,
CounterexamplesAbsent---are each independently necessary and jointly
sufficient for standing.
Their conjunction is the \emph{Bounded Receipted Chatman Equation}.

The phase change is this:

\begin{quote}
\itshape
From systems that execute assumed law\\
to systems that admit law before execution.
\end{quote}

Configuration is the first operational law.
Bounded admission is the mechanism.
$\qc = 1$ is the gate.

The future is not autonomous generation.
The future is \emph{bounded admission} of machine-generated artifacts.

That is the claim.
That is the thesis.

% ============================================================
\appendix

\section{Proof of Merge Idempotency}

\begin{proposition}
For a single-layer configuration ($n = 1$),
$\bigsqcup_\sigma \{x^{(1)}\} = x^{(1)}$.
\end{proposition}
\begin{proof}
$W(k) = 1$ for all $k \in \dom(x^{(1)})$, so $x^* = x^{(1)}$.
\end{proof}

\section{BLAKE3 Properties}

BLAKE3 is a cryptographic hash function with the following properties
used in this paper:
\begin{enumerate}
  \item \emph{Determinism}: equal inputs produce equal outputs.
  \item \emph{Collision resistance}: finding $a \neq b$ with
        $\mathrm{BLAKE3}(a) = \mathrm{BLAKE3}(b)$ requires $O(2^{128})$ operations.
  \item \emph{Preimage resistance}: given $h$, finding $a$ with
        $\mathrm{BLAKE3}(a) = h$ requires $O(2^{256})$ operations.
  \item \emph{Avalanche effect}: changing one bit in the input changes,
        on average, half the bits of the output.
\end{enumerate}

\section{Glossary of Notation}

\begin{tabular}{@{}ll@{}}
\toprule
Symbol & Meaning \\
\midrule
$\C$ & Configuration space \\
$\A$ & Admitted set ($\qc = 1$) \\
$\A_0$ & Feasible region (BoundSatisfied) \\
$\F_k$ & Value domain of field $k$ \\
$\Sigma$ & Schema (set of field names) \\
$\phi_k$ & Constraint function for field $k$ \\
$\sqcup_\sigma$ & Precedence join (merge) under ordering $\sigma$ \\
$W$ & Winner map \\
$\Lc$ & Admission lattice \\
$\mathbf{Adm}$ & Monoidal admission category \\
$\W$ & Witness algebra (BLAKE3 digests under XOR) \\
$\omega$ & Witness function \\
$\ce$ & Counterexample detector set \\
$c_i$ & Individual counterexample detector \\
$\qc$ & Standing bit (Bounded Receipted Chatman Equation) \\
$\tilde{q}$ & Smooth surrogate for $\qc$ \\
$\Lambda_k$ & Sensitivity of standing to field $k$ \\
$f$ & Fitness score \\
$\mu$ & DfCM measure over variant space \\
$\rho$ & Combinatorial coverage \\
$\alpha$ & ANDON gate function \\
\bottomrule
\end{tabular}

% ============================================================
\end{document}
